\chapter{LITERATURE REVIEW}
\label{ch:literature}
This chapter reviews prior work on guiding users toward objects of interest in immersive environments and on how the voice that delivers guidance can shape the user's experience.

\section{Guiding Users Toward Objects of Interest}
Object search in virtual and augmented reality has been studied extensively in head-mounted displays~\cite{gruenefeld2017eyesee360,trepkowski2019fov,chiossi2024searching,quinn2024review,kelley2025cueing}. This has utility in augmented reality, teleoperation, and simulation training~\cite{binetti2021outofview}. Researchers have investigated a variety of methods to point individuals toward objects in virtual reality, from visual~\cite{grogorick2017subtle,gruenefeld2017eyesee360} to auditory~\cite{marquardt2020guidance,tanaka2026coarse}, or a combination of both~\cite{binetti2021outofview,zhang2023seehear}.

When an object begins outside the field of view, assistance must first support attention orientation, directing the user toward a relevant region of the scene~\cite{sabouni2025guidance}. Once that region becomes visible, the user still needs to distinguish the intended object from nearby distractors.

Speech has been studied as one potential way to interact with three-dimensional space in this way since Bolt's seminal work \emph{Put-That-There}, which combined spoken commands with pointing to place and move objects in a spatial display~\cite{bolt1980putthatthere}. Advances in virtual and augmented reality have renewed this area of research~\cite{zhang2023seehear,kim2026speechtospatial,tanaka2026coarse}, with possibilities for engagement with the environment being essentially endless. In addition, audio-based methods specifically could improve quality of life for blind and low-vision people, for whom spoken spatial guidance can make locating objects easier, helping make their lives more independent~\cite{coughlan2020spatial,tanaka2026coarse}.

Research on auditory guidance supports using sound to help users locate objects. In head-mounted AR, Binetti et al. combined visual arrows with spatialized alarm sounds to guide users toward out-of-view targets~\cite{binetti2021outofview}. On narrow field-of-view AR displays, Marquardt et al. found audio-tactile guidance slower than a comparable visualization technique but equally accurate, with users exhibiting better situational awareness~\cite{marquardt2020guidance}. In addition, Zhang et al.'s AR book-search study used audio and visual hints to help identify a specific book on a shelf~\cite{zhang2023seehear}. In VR, Tanaka et al. supplemented continuous spatial audio with on-demand spoken directions such as ``up'' while blind participants reached for targets; this verbal mode produced the lowest error rates, fastest completion, lowest workload, and strongest preference~\cite{tanaka2026coarse}. We take heavy inspiration from this type of vocal guidance in our research.

Follow My Voice applies this principle, starting with coarser spoken directions that orient the participant toward a particular target region, after which more precise hints refine the search within it~\cite{kwok2022twostep,tanaka2026coarse}. The agent, which is gaze-contingent, knows the whereabouts of the target, while the participant initially does not know which object to select. This creates a dependency on the agent's gaze-based instructions to find it.

\section{Self-Similar Voice}
Recent research indicates that the voice delivering an instruction can influence how users experience an interaction as well as how they engage with a task. Listeners recognize personality in synthesized speech and respond more favorably to synthetic voices whose personality resembles their own~\cite{nass2001personality}. The user's own self-similar voice has also emerged as a persuasive medium: AI-generated encouragement in one's own voice was perceived as uniquely engaging and personal, motivating the user~\cite{fang2025emotional}. Furthermore, another pioneering study showed that a self-similar voice made a virtual collaborator more likable and believable, though also more eerie~\cite{guo2024doppelganger}. Instant voice-cloning methods have made it practical to synthesize a voice resembling a specific speaker~\cite{qin2023openvoice}. In Follow My Voice, self-similar synthesis creates a cloned voice resembling the participant, allowing the study to examine whether that resemblance changes the experience or effectiveness of receiving spatial instructions in an object search scenario, in terms of performance and perceived workload, eye-tracking indicators of workload and search efficiency, and engagement in the search task (Section~\ref{sec:research-questions}).
